\cotefile{157-5}
\scanpage{36}\mspage{23}

\unsure{suffisantes} : bien ordonné dans un cas, loc fini dans un autre. L'idée naturelle, c'est de \deleted{\unsure{prouver}} donner un énoncé commun, qui coiffe les deux à la fois, avec $I$ \unsure{artinien}, i.e. toute suite décroissante est stationnaire.

Ce qu'il faut, c'est une hyp. sur $I$ d'une part, sur $u_\bullet:X_\bullet\longrightarrow \unsure{Z|Y}_{\unsure{\bullet}}$ dans $M(I^{\unsure{\circ|o}})$ d'autre part, qui assure ceci

$(\unsure{\alpha|*})$\qquad $\forall\,i_0\in I$, et toute partie ouverte $J'$ de $J=I_{/i_0}$, la flèche $u_{J',J}$ dans le diagramme

\diagram{\begin{tikzcd}[ampersand replacement=\&,column sep=large]
 X_{J'} \arrow[r] \arrow[d,"u_{J'}"'] \arrow[dr,phantom,"\text{cocart}",pos=.54] \&
 X_J \arrow[d,dash,"\widetilde u_{J',J}"] \arrow[dd,bend left=45,"u_J"] \\
 Y_{J'} \arrow[r] \& Y_{J',J} \arrow[d,"u_{J',J}"'] \\
 \& Y_J
\end{tikzcd}}

est dans $C$.

Question : est-il exact qu'il suffit, pour la validité de \unsure{$\alpha$}, de supposer

\[
\begin{aligned}
&\text{a)}\quad u_\bullet\in \unsure{\mathbb C|C|\mathcal C}_I \quad\text{i.e.}\quad u_{J',J}\in C
 \quad \text{quand } J\deleted{\unsure{\in}}\unsure{=|\in}I_{/i_0},\\
&\qquad J'=J\setminus\{i_0\},\quad\text{et}
\end{aligned}
\]

b)\quad $I$ \uline{loc. fini} \deleted{\unsure{?}}

% archive sheet number: 35

